\documentclass[11pt]{article}
\usepackage[a4paper,margin=27mm]{geometry}
\usepackage{amsmath,amssymb,amsthm,hyperref}
\newtheorem{theorem}{Theorem}\newtheorem{proposition}{Proposition}
\newcommand{\Ad}{\operatorname{Ad}}\newcommand{\ad}{\operatorname{ad}}
\title{Why the squarefree prefix Gram matrix does not force factorial rank}
\author{Asymptotic method diagnostic; independently reviewed by all three collaborators}
\date{30 September 2026}
\begin{document}\maketitle
Normalize each letter's total mass to one. A positive weighted word
is viewed as an axis path with cumulative coordinates $X_i(t)$;
the coordinate $X_i$ increases only on the segments carrying letter
$i$. For a word $u$ on distinct letters, write $F_u(t)$ for the
corresponding prefix iterated integral. Put $F_{\varnothing}=1$.
The endpoint signature of a fair family has coefficient $1/k!$
on every word of $k$ distinct letters.

\begin{theorem}[A low-rank positive completion]\label{gram}
Fix a marked letter $c$ and an integer $r\ge1$. Assume fairness
through degree $2r+1$. Index prefix functions by all distinct-letter
words $u$ of length at most $r$ avoiding $c$, and consider their
Gram matrix in $L^2(dX_c)$.
Fairness and the shuffle identity determine, whenever the supports
of $u,v$ are disjoint,
\[
 G_{uv}:=\int F_u(t)F_v(t)\,dX_c(t)
       =\frac1{|u|!\,|v|!\,(|u|+|v|+1)}.       \tag{1}
\]
These specified entries admit a positive semidefinite completion
of rank at most $r+1$, regardless of the size of the alphabet.
Restricted to words of a single fixed degree, they admit a rank-one
positive semidefinite completion.
\end{theorem}
\begin{proof}
For disjoint supports, the shuffle product contains
$\binom{a+b}{a}$ distinct interleavings, where $a=|u|$, $b=|v|$.
Appending the marked letter makes every resulting endpoint word
have $a+b+1$ distinct letters. Therefore
\[
 \int F_uF_v\,dX_c
 =\sum_{w\in u\mathbin{\mathrm{shuffle}}v}\operatorname{Sig}_{wc}
 =\frac{\binom{a+b}{a}}{(a+b+1)!},
\]
which is (1). Here $u\mathbin{\mathrm{shuffle}}v$ simply denotes the set of
order-preserving interleavings.

Now set the right side of (1) as the value of \emph{every} matrix
entry, including pairs with overlapping supports. This is the Gram
matrix on $[0,1]$ of the functions
\[
 f_u(t)=t^{|u|}/|u|!.
\]
It is positive semidefinite and its functions span a space of
dimension at most $r+1$. At fixed degree all the functions coincide,
so that block has rank one. This completion is an algebraic
compatibility statement, not a claim that every finite fair word
has those particular repeated-letter moments.
\end{proof}

\paragraph{What disjoint polarization actually says.}
For two orderings $u,u'$ of one support, put $D=F_u-F_{u'}$.
Equation (1) says that $D$ is orthogonal to every disjoint-support
prefix function covered by the degree bound. Two such order contrasts
on disjoint supports are also orthogonal. Their squared norms, however,
involve repeated letters and are not specified by (1). In the
positive completion above every contrast is identically zero.
Thus disjoint polarization and Gram positivity alone do not force
even nonzero norms for these contrasts, much less factorially many
independent prefix functions. Adding the mean moments (empty word)
does not repair the argument: they are already included in the same
rank-$(r+1)$ completion.

\begin{proposition}[The continuous diagonal realization]\label{diagonal}
The completion in Theorem~\ref{gram} is realized by the continuous
positive path $X_1(t)=\cdots=X_n(t)=t$, $0\le t\le1$. This path is
exactly fair in every squarefree degree. Its all-time endpoint
insertion directions, projected through Lie degree $r\le n$, span
a space of dimension at most
\[
 n+(r-1)(n-1).
\]
\end{proposition}
\begin{proof}
Every prefix iterated integral of degree $a$ is $t^a/a!$, and the
endpoint signature is $\exp(H)$, where $H=\sum_iX_i$ is now the
sum of the formal Lie generators. The insertion directions are
\[
 \Ad_{\exp(tH)}X_i
   =\sum_{a=0}^{r-1}\frac{t^a}{a!}(\ad H)^aX_i
                       \quad\text{through degree }r.
\]
The degree-zero-in-$t$ vectors span at most $n$ dimensions.
For each $a\ge1$, the $n$ vectors $(\ad H)^aX_i$ have the relation
$\sum_i(\ad H)^aX_i=(\ad H)^aH=0$, so they span at most $n-1$
dimensions. Summing these bounds proves the assertion.
\end{proof}

This continuous example is \emph{not} a finite axis-word construction.
It shows precisely why a factorial-rank proof must exploit the
restriction that only one coordinate moves at a time. Endpoint
fairness, nonnegative coordinates, shuffle identities, and positivity
of prefix Gram matrices remain compatible with the much smaller
rank in Proposition~\ref{diagonal}.

\paragraph{Why deleting repeated letters is not a positive repair.}
One cannot simply set all repeated-letter products to zero inside a
positive moment argument. In the commutative reduced shuffle algebra,
a nonempty monomial $a$ has $a^2=0$. A unital positive functional
would have
\[
 0\le L((1+ta)^2)=1+2tL(a)\qquad(t\in\mathbb R),
\]
forcing $L(a)=0$. In particular this contradicts the nonzero
first-letter moments of the fair target. Equivalently, the proposed
moment matrix on $1,a$ would have a nonzero off-diagonal entry and
a zero diagonal entry, hence would fail positivity. Repeated-letter
entries must therefore be supplied rather than discarded; the
low-rank completion above is one legitimate supply at the level
of these moment constraints.

\paragraph{Finite words and the remaining question.}
There are separate genuine finite-word obstructions to full automatic
regularity. The reviewed note
\path{../../size_bounds/all_gap_singularity_inheritance.tex}
constructs arbitrarily large fully fair positive words whose all-gap
endpoint map is not surjective even through a fixed truncation degree
$r\ge5$. Its fixed singular seed is used only to falsify an asymptotic
universal lemma. The present argument is independent of that seed and
of its arithmetic certificate.

Neither result rules out a weaker factorial-order lower bound on the
rank of every finite fair word. Such a theorem would still be a
substantial route to a stronger total-size lower bound. It needs
additional information from the actual axis-word geometry or from
repeated-letter integrals; the squarefree Gram and disjoint-polarization
identities on their own cannot supply it. The separate standard-insertion
prefix collapse in
\path{../../independent_directions/prefix_signature_rank_obstruction.tex}
provides an exact finite-word warning against raw prefix independence.
\end{document}
